\documentclass[10pt,a4paper]{amsart}
\usepackage[margin=19mm]{geometry}
\usepackage{amsmath,amssymb,mathtools}
\usepackage[scr=rsfs,cal=euler]{mathalfa}
\usepackage{xfrac}
\usepackage[unicode,hidelinks,bookmarks=false]{hyperref}
\newcommand{\ie}{i.e.\ }
\newcommand{\cf}{cf.\ }
\newcommand{\resp}{respectively\ }
\newcommand{\Subsection}{\subsection}
\providecommand{\nobreakdash}{\nobreak\hbox{-}}
\setlength{\emergencystretch}{2em}
\multlinegap0pt
\usepackage{bm}
\usepackage{bbm}
\usepackage{dsfont}
\usepackage{xspace}
\usepackage{cancel}
\usepackage{mathtools}
\usepackage{amsthm}
\makeatletter
\@ifundefined{Theorem}{\newtheorem{Theorem}{Theorem}}{}
\@ifundefined{Lemma}{\newtheorem{Lemma}[Theorem]{Lemma}}{}
\@ifundefined{defprop}{\newtheorem{defprop}[Theorem]{Definition/Proposition}}{}
\makeatother
\DeclareMathOperator{\Ker}{Ker}
\DeclareMathOperator{\Ran}{Ran}
\DeclareMathOperator{\spann}{span}
\title[Perturbations of APS Boundary Conditions for Lorentzian Dira]{Perturbations of APS Boundary Conditions for Lorentzian Dirac Operators (Lemma 3.8)}
\author{Lennart Ronge}
\date{Source: arXiv:2408.16696v3 (CC BY 4.0)}
\begin{document}
\maketitle

\noindent\emph{Source and licence.} Perturbations of APS Boundary Conditions for Lorentzian Dirac Operators. Version arXiv:2408.16696v3 (CC BY 4.0), distributed under the Creative Commons Attribution 4.0 International licence, as recorded in the arXiv licence metadata for this exact version and shown on the abstract page https://arxiv.org/abs/2408.16696v3. Full rights notice: licenses/arxiv-2408.16696-CC-BY-4.0.txt.\par\smallskip

\noindent\textit{Editorial scope.} Counted unit: the Lemma printed as Lemma 3.8 of the article (source0.tex lines 422--427, source label \texttt{graphest}), delivered together with the article's own complete original proof of it (source0.tex lines 431--477, one \texttt{proof} environment of 47 lines). No mathematics is written, repaired or re-derived: statement and proof are the article's own text line for line. The proof is detached from the statement in the source: the article states the result at lines 422--427 and prints its proof at lines 431--477, and the proof environment's own header names the statement, so the association is the article's own. Required context retained uncounted: cnorm (lines 119--129); 2dE (lines 228--231). Every definition, display and label of those lines is retained unchanged. Omitted from this package and not redistributed: 1--118, 130--227, 232--421, 428--430, 478--642. Nothing inside a retained excerpt is omitted or altered: every retained body is delivered exactly as the article's lines are written, so no omission receipt of any kind accompanies this package. Citations: the retained excerpts carry no citation command at all, so no bibliography entry of the article is transcribed and no reference is redistributed. Reference adaptation: none was needed and none was made; every reference inside the delivered unit resolves inside it, so no unresolved reference or citation is produced. Printed slips kept literal: no printed word, symbol, hypothesis or number of the counted statement or of its proof was altered. Layout and class: the article's own class is \texttt{sigma} with packages including ifpdf, tikz; the wrapper is the lane's pinned \texttt{amsart} editorial wrapper at 10pt on A4, which restates the theorem-like environments the retained text needs and adds the article's own macro definitions that the retained text needs (\texttt{\textbackslash Ker}, \texttt{\textbackslash Ran}, \texttt{\textbackslash spann}), with extra wrapper packages (bm, bbm, dsfont, xspace, cancel, mathtools). The delivered numbering is the unit's own. Assets: no retained span contains a figure environment, a graphics-inclusion command, a PDF-page inclusion, a table or a TikZ picture; no image file is copied into this package, the package declares no asset dependency and no image file is redistributed with it. Licence: version arXiv:2408.16696v3 of this article is distributed under the Creative Commons Attribution 4.0 International licence, as recorded in the arXiv licence metadata for this exact version and shown on the abstract page https://arxiv.org/abs/2408.16696v3. A full rights notice is delivered at licenses/arxiv-2408.16696-CC-BY-4.0.txt. Characters: every retained excerpt is pure ASCII, so the delivered engine, XeLaTeX with its default Latin Modern text font, reports no missing character.
	\begin{defprop}		\label {cnorm}
		For $T\in\mathcal{B}(H)$, denote by $\|T\|_C$ the norm of the equivalence class of $T$ in $\mathcal{C}(H)$. This is equal to each of the following:
		\begin{enumerate}\itemsep=0pt
			\item[$(1)$] %\label{1}
            $\|T\|_1:=\inf\{\|S\| \mid S\in\mathcal{B}(H),\, S-T\in\mathcal{K}(H)\}$.
			\item[$(2)$] %\label{2}
            $\|T\|_2:=\inf\{\|S\| \mid S\in\mathcal{B}(H),\, S-T\text{ has finite rank}\}$.
			\item[$(3)$] %\label{3}
            $\|T\|_3:=\inf\{\|T|_V\|\mid V\subseteq H \text{ is a closed subspace of finite codimension}\}$.
		\end{enumerate}
	\end{defprop}

	\begin{Lemma}
		\label{2dE}
		Let $p$ be the projection onto $\spann\{(a,b)\}$ in $\mathbb{C}^2$, for $a,b\in \mathbb{R}$, and let $q$ be the projection onto $\mathbb{C}\times0$. Then $\|p-q\|=\frac{|b|}{\sqrt{a^2+b^2}}$.
	\end{Lemma}

	\begin{Lemma}
		\label{graphest}
		Let $H=Y\oplus Z$ be an orthogonal decomposition of $H$. Let $G\colon Y\rightarrow Z$ be a bounded linear map. Let $P$ be the projection onto the graph of~$G$, viewed as a subset of $H$, i.e.,
		\[\Gamma(G):=\{x+Gx\mid x\in Y\}.\]
		Let $S$ denote the projection onto $Y$. Then \smash{$\|P-S\|_C\leq\frac{\|G\|_C}{\sqrt{1+\|G\|_C^2}}$}.
	\end{Lemma}

	\begin{proof}[Proof of Lemma \ref{graphest}]
		Before we show the estimate, there are many simplifications we can do without loss of generality.

		Firstly, we want to replace Calkin norms with operator norms. Suppose we had only shown the theorem with $\|G\|$ instead of $\|G\|_C$ on the right-hand side. Let $\varepsilon >0$ and let $Y'\subset Y$ be a~subspace of finite codimension such that
		\[\|G|_{Y'}\|\leq\|G\|_C+\varepsilon\]
		(this exists by Definition/Proposition~\ref{cnorm}). Then applying the weaker version of the lemma with~$Y'$ replacing $Y$ and $G':=G|_{Y'}$, we can deduce that
		\[\bigl\|P'-S'\bigr\|_C\leq \frac{\|G'\|}{\sqrt{1+\|G'\|^2}},\]
		where $S'$ is the projection onto $Y'$ and $P'$ that on the graph of $G'$ in $Y'\oplus Z$. Extending them by zero \big(on $Y'^\bot\subset Y$\big), we obtain the projections on these spaces in $H$, which have compact difference from $S$ (resp.\ $P$) (the ranges are finite codimension subsets of each other).
		Thus
		\[\|P-S\|_C=\bigl\|\bigl(P'-S'\bigr)\oplus 0\bigr\|_C=\bigl\|P'-S'\bigr\|_C\leq \frac{\|G'\|}{\sqrt{1+\|G'\|^2}}\leq \frac{(\|G\|_C+\varepsilon)}{\sqrt{1+(\|G\|_C+\varepsilon)^2}}.\]
		As $\varepsilon$ was arbitrary, this implies the stronger lemma. Thus it suffices to show the lemma without the subscript $C$ on the right-hand side. As we always have $\|P-S\|_C\leq \|P-S\|$, it is sufficient to show the lemma with operator norms everywhere.
		
		Secondly, we may assume that the kernel of $G$ is zero and it has dense range: As both $P$ and~$S$ are~$1$ on $\Ker(G)$ and $0$ on $\Ran(G)^\bot$, their difference is zero on either space. Thus removing these spaces does not change the left-hand side of the inequality. As it also does not change $\|G\|$ and hence the right-hand side, it suffices to show the theorem in the case where
\[
\Ker(G)=\Ran(G)^\bot=0,
\]
		which we will assume from now on.

		Thirdly, we want to replace $G$ by a self-adjoint operator. Let $G=UA$ be the polar decomposition of $G$, i.e., $A$ is self-adjoint and $U$ is an isometry from $\Ker(G)^\bot$ to $\overline{\Ran(G)}$. By the previous assumption, these spaces are everything, i.e., $U$ is a unitary map from $Y$ to $Z$. Suppose we had shown the lemma for $A$ instead of $G$ (and $Y$ instead of $Z$) and let again $P'$ and $S'$ be the projections for that case. As the graph of $G$ is the image of the graph of $A$ under $\tilde U:=1\oplus U$ (which also fixes $Y\oplus 0$), we get
		\[
        \|P-S\|=\bigl\|\tilde U\bigl(P'-S'\bigr)\tilde U^{-1}\bigr\|=\bigl\|P'-S'\bigr\|\leq \frac{\|A\|}{\sqrt{1+\|A\|^2}}=\frac{\|G\|}{\sqrt{1+\|G\|^2}}.
        \]
		Thus we may assume from now on that $Z=Y$ and $G$ is self-adjoint.

		Fourthly, we want to replace $G$ by a multiplication operator. By the spectral theorem, there is a unitary $V\colon Y\to L^2(\Omega)$ for some measure space $\Omega$ and an $L^\infty$-function $g$ such that $VGV^{-1}f=gf$ for any $f\in L^2(\Omega)$ (in particular $\|G\|=\|g\|_\infty$). As everything is invariant under conjugation by unitaries, it suffices to show the lemma for $Y=L^2(\Omega)$ and $G$ a multiplication operator.
		
		Finally, we identify $L^2(\Omega)\oplus L^2(\Omega)$ with $L^2\bigl(\Omega,\mathbb{C}^2\bigr)$.
		
		Thus, overall, have the following situation: Let $\Omega$ be a measure space with measure $\mu$ and $g\in L^\infty(\Omega)$. In $L^2\bigl(\Omega, \mathbb{C}^2\bigr)$, let $S$ be the projection onto $L^2(\Omega,\mathbb{C}\times\{0\})$ and $P$ the projection onto $\bigl\{(\psi,g\psi)\mid \psi\in L^2(\Omega)\bigr\}$. We need to show the inequality
		\[
        \|P-S\|\leq \frac{\|g\|_\infty}{\sqrt{1+\|g\|_\infty^2}}.
        \]
		Denote by $p_a$ the projection onto $\spann\{(1,a)\}$ in $\mathbb{C}^2$. Then the following define orthogonal projections onto the correct spaces:
		\begin{gather*}
        Sf(x)=p_0f(x),\qquad
        Pf(x)=p_{g(x)}f(x).
        \end{gather*}
		We can thus check for any $f\in H$, using Lemma \ref{2dE} for the pointwise estimate:
		\begin{align*}
			\|(P-S)f\|^2
			&=\int\limits_\Omega\|(p_{g(x)}-p_0)f(x)\|^2{\rm d}\mu(x)
			 \leq\int\limits_\Omega\frac{g(x)^2}{1+g(x)^2}\|f(x)\|^2{\rm d}\mu(x)\\
			&\leq\frac{\|g\|_\infty^2}{1+\|g\|_\infty^2}\int\limits_\Omega\|f(x)\|^2{\rm d}\mu(x)
			 =\frac{\|g\|_\infty^2}{1+\|g\|_\infty^2}\|f\|^2,
		\end{align*}
		which gives the desired estimate.
	\end{proof}

\end{document}
